\documentclass[12pt, a4paper]{article}

%% ─── Packages ───────────────────────────────────────────────────────────────
\usepackage[a4paper, margin=2.5cm]{geometry}
\usepackage{amsmath, amssymb}
\usepackage{graphicx}
\usepackage{booktabs}
\usepackage{array}
\usepackage{caption}
\usepackage{subcaption}
\usepackage{hyperref}
\usepackage{fancyhdr}
\usepackage{titlesec}
\usepackage{xcolor}
\usepackage{enumitem}
\usepackage{siunitx}
\usepackage{float}
\usepackage{listings}
\usepackage{parskip}

\setlength{\headheight}{15pt}
\setlength{\parskip}{1em}

%% ─── Page Style ─────────────────────────────────────────────────────────────
\pagestyle{fancy}
\fancyhf{}
\rhead{EE568 -- Project 2}
\lhead{Elif Topaloğlu}
\cfoot{\thepage}
\renewcommand{\headrulewidth}{0.4pt}

%% ─── Hyperref Setup ─────────────────────────────────────────────────────────
\hypersetup{
    colorlinks=true,
    linkcolor=blue!60!black,
    urlcolor=blue!60!black,
    citecolor=blue!60!black,
}

%% ─── Section Formatting ─────────────────────────────────────────────────────
\titleformat{\section}{\large\bfseries}{\thesection\quad}{0em}{}
\titleformat{\subsection}{\normalsize\bfseries}{\thesubsection\quad}{0em}{}

%% ─── SI Units Setup ─────────────────────────────────────────────────────────
\sisetup{per-mode=symbol}

%% ─── MATLAB Listing Style ───────────────────────────────────────────────────
\lstdefinestyle{matlab}{
    language=Matlab,
    basicstyle=\ttfamily\small,
    keywordstyle=\color{blue}\bfseries,
    commentstyle=\color{green!50!black},
    stringstyle=\color{orange},
    numbers=left,
    numberstyle=\tiny\color{gray},
    numbersep=8pt,
    frame=single,
    breaklines=true,
    captionpos=b,
    tabsize=4,
}

\begin{document}

%%─── Title Page ──────────────────────────────────────────────────────────────
\begin{titlepage}
    \centering
    \begin{minipage}{0.25\textwidth}
        \centering
        
\includegraphics[height=3cm]{figures/eelogo.png} 
    \end{minipage}
    \hfill
    \begin{minipage}{0.25\textwidth}
        \centering
        
\includegraphics[height=3cm]{figures/Orta_Doğu_Teknik_Üniversitesi.png}
    \end{minipage}
    
    \vspace{2cm}
    {\LARGE \textbf{Middle East Technical University} \par}
    \vspace{1cm}
    {\large Department of Electrical and Electronics Engineering \par}
    \vspace{2cm}
    \rule{\linewidth}{0.5mm}\\[0.5cm]
    {\huge \bfseries EE568 -- Selected Topics in Electrical Machines \par}
    {\LARGE \bfseries Project 2: Motor Winding Design \& Analysis \par}
    \rule{\linewidth}{0.5mm}\\[1.5cm]
    
    \begin{tabular}{rl}
        \textbf{Author:}    & Elif Topaloğlu \\[0.5em]
        \textbf{Date:}      & \today \\[0.5em]
    \end{tabular}
    \vfill
    {\large \textbf{GitHub Repository:} \\[0.3em]
    \url{https://github.com/odtu/EE568/tree/master/Project2}}
\end{titlepage}
%%──────────────────────────────────────────────────────────





%%─── Table of Contents ───────────────────────────────────────────────────────
\tableofcontents
\newpage

%%─── 1. Introduction ─────────────────────────────────────────────────────────
\section{Introduction}

This report mainly presents the winding design and analytical modelling work carried out for EE568 Project 2. The report is organised into four main parts. In the first part, an integral-slot double-layer winding is designed for a 6-pole, 72-slot, 3-phase machine. Both full-pitched and 11/12 short-pitched configurations are analysed, and the distribution, pitch, and winding factors are calculated for the fundamental, 3rd, and 5th harmonics. The effect of short-pitching on harmonic suppression is discussed.

In the second part, a fractional-slot concentrated winding is designed for the assigned 15-slot/14-pole permanent-magnet synchronous machine. The induced EMF phasor diagram is constructed for each slot, the winding factor is computed for the fundamental and selected harmonics, and the balance of the three-phase winding is verified analytically.

The third part presents the analytical modelling of the machine, covering the electrical loading, Carter's coefficient, effective air-gap length, air-gap flux density, magnetic loading, specific machine constant, shear stress, expected
torque, and flux per pole.

Finally, in the fourth part, a 2D finite element analysis (FEA) is performed using ANSYS Maxwell to verify the analytical results. The flux density distribution, air-gap flux density, leakage flux paths, flux per pole, etc.are obtained and discussed.



\section{Q1: Integral-Slot Winding Design}

For this question, two winding designs are done namely full-pitched and 11/12 short-pitched windings. The given machine parameters for these windings are:

\begin{itemize}
    \item Number of poles: $p = 6$ 
    \item Number of stator slots: $Q = 72$
    \item Number of phases: $m = 3$
    \item Winding configuration: Double-layer
\end{itemize}

\paragraph{Derived Parameters:}

\begin{enumerate}
    \item \textbf{Slots per pole per phase ($q$):}
    \[ q = \frac{Q}{p \cdot m} = \frac{72}{6 \times 3} = 4\]

    \item \textbf{Electrical slot angle ($\alpha_e$):}
    \[\alpha_e = \frac{p \times 180^\circ}{Q} = \frac{6 \times 180^\circ}{72} = 15^\circ/\text{slot}\]

    \item \textbf{Mechanical slot angle ($\alpha_m$):}
    \[\alpha_m = \frac{360^\circ}{Q} = \frac{360^\circ}{72} = 5^\circ/\text{slot}\]

    \item \textbf{Pole pitch ($\tau_p$):}
    \[\tau_p = \frac{Q}{p} = \frac{72}{6} = 12 \text{ slots}\]
\end{enumerate}


\subsection{Configuration A: Full-Pitched Winding}

\subsubsection{Winding Diagram}

For a full-pitched winding, the coil span equals the pole pitch: $y = \tau_p = 12$ slots, corresponding to $180^\circ$ electrical. In a double-layer winding, the top layer conductor in slot $i$ is connected to the bottom layer conductor in slot $i + y$. This means the bottom layer of slot $i$ serves as the return side of the coil whose top side resides in slot $i - y$. Since the phase sequence repeats every 12 slots with alternating polarity, the bottom layer ends up identical to the top layer at each slot position.

The phase sequence follows the order $A,\; {-C},\; B,\; {-A},\; C,\; {-B}$, giving a balanced 3-phase layout across the 24-slot pole-pair:

\begin{table}[ht!]
\centering
\footnotesize
\setlength{\tabcolsep}{2.5pt}
\caption{Double-layer full-pitched winding layout (one pole-pair)}
\begin{tabular}{l|cccc|cccc|cccc|cccc|cccc|cccc}
\toprule
\textbf{Slot} & 1 & 2 & 3 & 4 & 5 & 6 & 7 & 8 & 9 & 10 & 11 & 12 & 13 & 14 & 15 & 16 & 17 & 18 & 19 & 20 & 21 & 22 & 23 & 24 \\
\midrule
\textbf{Top} & A & A & A & A & -C & -C & -C & -C & B & B & B & B & -A & -A & -A  & -A  & C & C & C & C & -B & -B & -B & -B \\
\textbf{Bottom} & A & A & A & A & -C & -C & -C & -C & B & B & B & B & -A & -A & -A  & -A  & C & C & C & C & -B & -B & -B & -B \\
\bottomrule
\end{tabular}
\end{table}

\subsubsection{Winding Factor Calculations}

The general formulas for the $n$-th harmonic are:
\[
    k_{d,n} = \frac{\sin\!\left(\dfrac{n \cdot q \cdot \alpha_e}{2}\right)}{q \cdot \sin\!\left(\dfrac{n \cdot \alpha_e}{2}\right)}
    \qquad
    k_{p,n} = \sin\!\left(\frac{n \cdot y_e}{2}\right)
    \qquad
    k_{w,n} = k_{d,n} \cdot k_{p,n}
\]

where $y_e = 180^\circ$ electrical for full pitch.

\paragraph{Fundamental ($n = 1$):}
\[
    k_{d,1} = \frac{\sin(4 \times 15^\circ / 2 \times 1)}{4\,\sin(15^\circ/2 \times 1)}
            = \frac{\sin(30^\circ)}{4\,\sin(7.5^\circ)}
            = \frac{0.500}{4 \times 0.131} = \frac{0.500}{0.522} = 0.958
\]
\[
    k_{p,1} = \sin(90^\circ) = 1.000 \qquad \Rightarrow \qquad k_{w,1} = 0.958
\]

\paragraph{3rd Harmonic ($n = 3$):}
\[
    k_{d,3} = \frac{\sin(3 \times 30^\circ)}{4\,\sin(3 \times 7.5^\circ)}
            = \frac{\sin(90^\circ)}{4\,\sin(22.5^\circ)}
            = \frac{1.000}{4 \times 0.383} = \frac{1.000}{1.531} = 0.653
\]
\[
    k_{p,3} = \sin(3 \times 90^\circ) = \sin(270^\circ) = -1.000 \qquad \Rightarrow \qquad k_{w,3} = -0.653
\]

\paragraph{5th Harmonic ($n = 5$):}
\[
    k_{d,5} = \frac{\sin(5 \times 30^\circ)}{4\,\sin(5 \times 7.5^\circ)}
            = \frac{\sin(150^\circ)}{4\,\sin(37.5^\circ)}
            = \frac{0.500}{4 \times 0.609} = \frac{0.500}{2.435} = 0.205
\]
\[
    k_{p,5} = \sin(5 \times 90^\circ) = \sin(450^\circ) = 1.000 \qquad \Rightarrow \qquad k_{w,5} = 0.205
\]

\begin{table}[h!]
\centering
\caption{Winding factors for the full-pitched configuration.}
\begin{tabular}{cccc}
\toprule
\textbf{Harmonic ($n$)} & \textbf{$k_{d,n}$} & \textbf{$k_{p,n}$} & \textbf{$k_{w,n}$} \\
\midrule
1 (Fundamental) & $0.958$ & $1.000$ & $0.958$ \\
3               & $0.653$ & $-1.000$ & $-0.653$ \\
5               & $0.205$ & $1.000$ & $0.205$ \\
\bottomrule
\end{tabular}
\end{table}


\subsection{Configuration B: 11/12 Short-Pitched Winding}

\subsubsection{ Winding Diagram}

For the $11/12$ short-pitched design, the coil span is reduced by one slot: $y = 11$ slots. In electrical degrees:
\[
    y_e = 11 \times 15^\circ = 165^\circ \text{ electrical}
\]
The top layer is identical to the full-pitch case. However, the bottom layer of slot $i$ now carries the return conductor of the coil whose top side is at slot $i - 11$. This shifts the bottom layer by one slot relative to the full-pitch arrangement:

\begin{table}[h!]
\centering
\footnotesize
\setlength{\tabcolsep}{2pt}
\caption{Double-layer 11/12 short-pitched winding layout.}
\begin{tabular}{l|cccc|cccc|cccc|cccc|cccc|cccc}
\toprule
\textbf{Slot} & 1 & 2 & 3 & 4 & 5 & 6 & 7 & 8 & 9 & 10 & 11 & 12 & 13 & 14 & 15 & 16 & 17 & 18 & 19 & 20 & 21 & 22 & 23 & 24 \\
\midrule
\textbf{Top}& A & A & A & A & -C & -C & -C & -C & B & B & B & B & -A & -A & -A  & -A  & C & C & C & C & -B & -B & -B & -B \\
\textbf{Bottom} & A & A & A & -C & -C & -C & -C & B & B & B & B & -A & -A & -A & -A  & C & C & C & C & -B & -B & -B & -B & A \\
\bottomrule
\end{tabular}
\end{table}

\subsubsection{Winding Factor Calculations}
The distribution factor does not depend on the pitch which makes the distribution factor identical for both configurations. Only the pitch factor changes:

\[k_{p,n} = \sin\!\left(\frac{n \times 165^\circ}{2}\right) = \sin(n \times 82.5^\circ)\]

\paragraph{Fundamental ($n = 1$):}
\[k_{p,1} = \sin(82.5^\circ) = 0.991 \qquad \Rightarrow \qquad k_{w,1} = 0.958 \times 0.991 = 0.949\]

\paragraph{3rd Harmonic ($n = 3$):}
\[k_{p,3} = \sin(247.5^\circ) = -0.924 \qquad \Rightarrow \qquad k_{w,3} = 0.653 \times (-0.924) = -0.604\]

\paragraph{5th Harmonic ($n = 5$):}
\[k_{p,5} = \sin(412.5^\circ) = \sin(52.5^\circ) = 0.793 \qquad \Rightarrow \qquad k_{w,5} = 0.205 \times 0.793 = 0.163\]

\begin{table}[h!]
\centering
\caption{Winding factors for the 11/12 short-pitched configuration.}
\begin{tabular}{cccc}
\toprule
\textbf{Harmonic ($n$)} & \textbf{$k_{d,n}$} & \textbf{$k_{p,n}$} & \textbf{$k_{w,n}$} \\
\midrule
1 (Fundamental) & $0.958$ & $0.991$ & $0.949$ \\
3               & $0.653$ & $-0.924$ & $-0.604$ \\
5               & $0.205$ & $0.793$ & $0.163$ \\
\bottomrule
\end{tabular}
\end{table}

\subsection{Discussion}
The following conclusions can be drawn:

\begin{enumerate}
    \item Moving from full-pitch to 11/12 short-pitch reduces the fundamental winding factor minimally, from $k_{w,1} = 0.958$ to $0.949$ (a decrease of 0.9\%). The machine sustains nearly the same torque and back-EMF capability.

    \item Third harmonic slightly decreased (7.5\%) for short-pitch configuration. However, in a balanced 3-phase machine the triplen harmonics ($n = 3, 9, 15, \ldots$) cancel in the \emph{line-to-line} voltage and do not appear in the output regardless of their winding factor. Yet, they may still contribute to stator core losses and noise.

    \item Short-pitching reduces $|k_{w,5}|$ from $0.205$ to $0.163$ ( approximately 20\% reduction). The pitch factor $k_{p,5}$ drops from $1.000$ to $0.793$. This is the primary benefit of the 11/12 pitch: it meaningfully attenuates the 5th space harmonic, which is responsible for additional losses in the machine.

    \item Short-pitching always involves a trade-off between harmonic suppression and fundamental EMF and torque. In this case the 11/12 pitch is a very efficient compromise, it achieves noticeable 5th harmonic reduction at the cost of less than 1\% in fundamental performance.
\end{enumerate}

\begin{table}[h!]
\centering
\caption{Winding factor comparison between the two configurations.}
\begin{tabular}{lcccc}
\toprule
 & \multicolumn{2}{c}{\textbf{Full Pitch ($y = 12$)}} & \multicolumn{2}{c}{\textbf{Short Pitch ($y = 11$)}} \\
\cmidrule(lr){2-3}\cmidrule(lr){4-5}
\textbf{Harmonic} & $k_{p,n}$ & $k_{w,n}$ & $k_{p,n}$ & $k_{w,n}$ \\
\midrule
$n=1$ & $1.000$ & $0.958$ & $0.991$ & $0.949$ \\
$n=3$ & $-1.000$ & $-0.653$ & $-0.92$ & $-0.604$ \\
$n=5$ & $1.000$ & $0.205$ & $0.793$ & $0.163$ \\
\bottomrule
\end{tabular}
\end{table}



\section{Q2: Fractional-Slot Winding Design}
The assigned design for this project has the following parameters:

\begin{itemize}
    \item Number of stator slots: $N_s = 15$
    \item Number of rotor poles: $N_m = 14$
    \item Number of phases: $m = 3$
    \item Slots per pole per phase: $N_{spp} \approx 0.357$
\end{itemize}

\subsection{Phase Angles of Induced Voltage per Slot}

The electrical angle of a slot can be found as:
\[\varphi_k = \frac{N_m \cdot 180^\circ}{N_s} \cdot k = \frac{14 \cdot 180^\circ}{15} = 168^\circ\ \cdot k\]
where $k$ is the slot index starting from 1. This means each slot's induced EMF will have a phase angle that is a multiple of $168^\circ$. Since $168^\circ$ is not an integer divisor of $360^\circ$, the slot phasors will be distributed around the circle in a non-repetitive manner, which is characteristic of fractional-slot windings. The angle of the induced EMF in every slot is presented in the Table~\ref{tab:induced_emf_1} below. Also, phaseor diagram of all slot is presented in Figure~\ref{fig:phasor_slots}.

\begin{table}[h!]
\centering
\caption{Induced EMF phase angle and phase assignment for each slot (fundamental, $n=1$).}
\label{tab:induced_emf_1}
\begin{tabular}{cccc|cccc}
\toprule
\textbf{Slot} & $\varphi_k$ & \textbf{Phase} & \textbf{Sign} &
\textbf{Slot} & $\varphi_k$ & \textbf{Phase} & \textbf{Sign} \\
\midrule
 1 & $0^\circ$ & A & $+$ &  9 & $264^\circ$ & C & $+$ \\
 2 & $168^\circ$ & A & $-$ & 10 &  $72^\circ$ & C & $-$ \\
 3 & $336^\circ$ & A & $+$ & 11 & $240^\circ$ & C & $+$ \\
 4 & $144^\circ$ & B & $+$ & 12 &  $48^\circ$ & C & $-$ \\
 5 & $312^\circ$ & B & $-$ & 13 & $216^\circ$ & C & $+$ \\
 6 & $120^\circ$ & B & $+$ & 14 &  $24^\circ$ & A & $+$ \\
 7 & $288^\circ$ & B & $-$ & 15 & $192^\circ$ & A & $-$ \\
 8 &  $96^\circ$ & B & $+$ &    &        &   &     \\
\bottomrule
\end{tabular}
\end{table}

\begin{figure}[H]
    \centering
    \includegraphics[width=0.7\textwidth]{figures/phasor_slots.eps}
    \caption{Phasor diagram of the fractional-slot winding.}
    \label{fig:phasor_slots}
\end{figure}

\subsection{Fundamental Harmonic (\texorpdfstring{$n=1$}{n=1})}

Phasor diagram for the fundamental harmonic ($n=1$) is shown in Figure~\ref{fig:phasor_three_phases}. The slot phasors for each phase are summed vectorially to find the resultant phasor for that phase. The angles of the resultant phasors for the three phases are $0^\circ$ for Phase A, $120^\circ$ for Phase B, and $240^\circ$ for Phase C, confirming a balanced three-phase system.

\begin{figure}[H]
    \centering
    \includegraphics[width=0.7\textwidth]{figures/phasor_three_phases.eps}
    \caption{Phasor diagram of the three phases ($n=1$).}
    \label{fig:phasor_three_phases}
\end{figure}

The winding factor for the $n$-th harmonic is computed as the ratio of the vector sum magnitude to the arithmetic sum magnitude of the slot phasors for one phase.
\[k_{w,n} = \frac{\left|\displaystyle\sum_{\text{Phase}} \vec{e}_{k,n}\right|}{N_s / m}\]
where each phasor $\vec{e}_{k,n}$ has unit magnitude. The Phase A phasors and their effective angles are presented in the Table~\ref{tab:phasors_phase_a1} below. Distribution factor for the first fundamental is calculated using the angles of the phasors.

\begin{table}[h!]
\centering
\caption{Phase A phasors for $n=1$.}
\label{tab:phasors_phase_a1}
\begin{tabular}{cccc}
\toprule
\textbf{Slot} & \textbf{Sign} & \textbf{Raw angle $\varphi_k$} & \textbf{Effective phasor angle} \\
\midrule
 1 & $+$ & $  0^\circ$ & $  0^\circ$ \\
14 & $+$ & $ 24^\circ$ & $ 24^\circ$ \\
15 & $-$ & $192^\circ$ & $ 12^\circ$ \\
 2 & $-$ & $168^\circ$ & $348^\circ$ \\
 3 & $+$ & $336^\circ$ & $336^\circ$ \\
\bottomrule
\end{tabular}
\end{table}

The vector sum of these 5 unit phasors is calculated by summing their real and imaginary components:
\[\sum \cos\theta = \cos(0^\circ) + \cos(24^\circ) + \cos(12^\circ) + \cos(348^\circ) + \cos(336^\circ) = 4.783\]
\[\sum \sin\theta = \sin(0^\circ) + \sin(24^\circ) + \sin(12^\circ) + \sin(348^\circ) + \sin(336^\circ) = 0\]
\[k_{w,1} = \frac{|\,\vec{e}_1 + \vec{e}_{14} + \vec{e}_{15} + \vec{e}_2 + \vec{e}_3\,|}{5} = \frac{4.783}{5} = \mathbf{0.957}\]

\subsection{Third Harmonic (\texorpdfstring{$n=3$}{n=3})}

For the 3rd harmonic, the electrical angle of each slot scales by a factor of 3. The effective electrical slot pitch becomes:
\[\alpha_{e,3} = 3 \times 168^\circ = 504^\circ \equiv 144^\circ\]

Phasor diagram for the third harmonic is shown in Figure~\ref{fig:phasor3}. The phasors for the 3rd harmonic will be more widely spaced due to the increased electrical angle, and their vector sum will yield a smaller resultant, indicating a lower winding factor for this harmonic.

\begin{figure}[H]
    \centering
    \includegraphics[width=0.7\textwidth]{figures/phasor3.eps}
    \caption{Phasor diagram of phase A ($n=3$).}
    \label{fig:phasor3}
\end{figure}

Multiplying the fundamental raw angles of Phase A by 3 and applying a $180^\circ$ shift for the negative-polarity coils yields the 3rd harmonic effective phasor angles shown in Table~\ref{tab:phasors_phase_a3}.

\begin{table}[h!]
\centering
\caption{Phase A phasors for $n=3$.}
\label{tab:phasors_phase_a3}
\begin{tabular}{cccc}
\toprule
\textbf{Slot} & \textbf{Sign} & \textbf{Raw angle $\varphi_k$} & \textbf{Effective Phasor Angle} \\
\midrule
 1  & $+$ & $0^\circ$   & $0^\circ$ \\
 14 & $+$ & $72^\circ$  & $72^\circ$ \\
 15 & $-$ & $216^\circ$ & $36^\circ$ \\
 2  & $-$ & $144^\circ$ & $324^\circ$ \\
 3  & $+$ & $288^\circ$ & $288^\circ$ \\
\bottomrule
\end{tabular}
\end{table}

The vector sum of these 5 unit phasors is calculated by summing their real and imaginary components:
\[\sum \cos\theta = \cos(0^\circ) + \cos(72^\circ) + \cos(36^\circ) + \cos(324^\circ) + \cos(288^\circ) = 3.236\]
\[\sum \sin\theta = \sin(0^\circ) + \sin(72^\circ) + \sin(36^\circ) + \sin(324^\circ) + \sin(288^\circ) = 0\]

The winding factor for the 3rd harmonic is computed as:
\[k_{w,3}  = \frac{3.236}{5} = \mathbf{0.647}\]


\subsection{Fifth Harmonic (\texorpdfstring{$n=5$}{n=5})}

For the 5th harmonic, the effective electrical slot pitch scales by a factor of 5:
\[\alpha_{e,5} = 5 \times 168^\circ = 840^\circ \equiv 120^\circ\]

Phasor diagram for the fifth harmonic is shown in Figure~\ref{fig:phasor5}.

\begin{figure}[H]
    \centering
    \includegraphics[width=0.7\textwidth]{figures/phasor5.eps}
    \caption{Phasor diagram of phase A ($n=5$).}
    \label{fig:phasor5}
\end{figure}

Multiplying the fundamental raw angles of Phase A by 5 and handling the coil polarities yields the 5th harmonic effective phasor angles presented in Table~\ref{tab:phasors_phase_a5}.

\begin{table}[h!]
\centering
\caption{Phase A phasors for $n=5$.}
\label{tab:phasors_phase_a5}
\begin{tabular}{cccc}
\toprule
\textbf{Slot} & \textbf{Sign} & \textbf{Raw angle $\varphi_k$} & \textbf{Effective Phasor Angle} \\
\midrule
 1  & $+$ & $0^\circ$   & $0^\circ$ \\
 14 & $+$ & $120^\circ$ & $120^\circ$ \\
 15 & $-$ & $240^\circ$ & $60^\circ$ \\
 2  & $-$ & $120^\circ$ & $300^\circ$ \\
 3  & $+$ & $240^\circ$ & $240^\circ$ \\
\bottomrule
\end{tabular}
\end{table}

The vector components sum up as follows:
\[\sum \cos\theta = \cos(0^\circ) + \cos(120^\circ) + \cos(60^\circ) + \cos(300^\circ) + \cos(240^\circ) = 1\]
\[\sum \sin\theta = \sin(0^\circ) + \sin(120^\circ) + \sin(60^\circ) + \sin(300^\circ) + \sin(240^\circ) = 0\]

The winding factor for the 5th harmonic is:
\[k_{w,5} = \frac{1}{5} = \mathbf{0.200}\]


\begin{table}[h!]
\centering
\caption{Winding factors summary for harmonics $n=1,3,5$.}
\begin{tabular}{ccc}
\toprule
\textbf{Harmonic ($n$)} & \textbf{$k_{w,n}$}\\
\midrule
$1$  & $0.957$ \\
$3$  & $0.647$ \\
$5$  & $0.200$ \\
\bottomrule
\end{tabular}
\end{table}


\subsection{Winding Diagram and Balance Verification}

The winding layout is designed as a double-layer concentrated configuration. The full physical slot distribution mapping all three phases across the 15 stator slots is presented in Table~\ref{tab:winding_layout_double}.

\begin{table}[h!]
\centering
\caption{Double-layer concentrated winding slot assignment mapping for all phases.}
\label{tab:winding_layout_double}
\resizebox{\textwidth}{!}{%
\begin{tabular}{c|ccccccccccccccc}
\toprule
\textbf{Slot ID} 
& 1 & 2 & 3 & 4 & 5 & 6 & 7 & 8 & 9 & 10 & 11 & 12 & 13 & 14 & 15 \\
\midrule

\textbf{Top Layer}
& A$+$ & A$-$ & A$+$ & B$+$ & B$-$ & B$+$ & B$-$ & B$+$ & C$+$ 
& C$-$ & C$+$ & C$-$ & C$+$ & A$+$ & A$-$ \\

\textbf{Bottom Layer}
& A$+$ & A$-$ & B$-$ & B$+$ & B$-$ & B$+$ & B$-$ & C$-$ & C$+$ 
& C$-$ & C$+$ & C$-$ & A$-$ & A$+$ & A$-$ \\

\bottomrule
\end{tabular}%
}
\end{table}


To verify that this fractional-slot double-layer layout represents a valid and balanced configuration, the net complex fundamental EMF voltage vectors ($\vec{E}_k = \text{Sign} \times e^{j \varphi_k}$) are computed for each phase using the baseline slot angle formula $\varphi_k = (k-1) \times 168^\circ$:

\begin{align*}
    \vec{E}_A &= 2e^{j0^\circ} - 2e^{j168^\circ} - 1e^{j216^\circ} + 2e^{j24^\circ} - 2e^{j192^\circ} = \mathbf{8.662 \angle 9.3^\circ} \\
    \vec{E}_B &= -1e^{j336^\circ} + 2e^{j144^\circ} - 2e^{j312^\circ} + 2e^{j120^\circ} - 2e^{j288^\circ} + 1e^{j96^\circ} = \mathbf{8.662 \angle 129.3^\circ} \\
    \vec{E}_C &= -1e^{j96^\circ} + 2e^{j264^\circ} - 2e^{j72^\circ} + 2e^{j240^\circ} - 2e^{j48^\circ} + 1e^{j216^\circ} = \mathbf{8.662 \angle 249.3^\circ}
\end{align*}

Evaluating the relative displacement between adjacent phases yields:
\begin{align*}
    \Delta\theta_{AB} &= 129.3^\circ - 9.3^\circ = 120.0^\circ \\
    \Delta\theta_{BC} &= 249.3^\circ - 129.3^\circ = 120.0^\circ \\
    \Delta\theta_{CA} &= 369.3^\circ - 249.3^\circ = 120.0^\circ
\end{align*}

The phase resultants form a balanced three-phase set. They possess completely uniform vector magnitudes and maintain an exact angular displacement of $120.0^\circ$ electrical degrees, confirming the electrical balance of the double-layer design.



\section{Q3: Analytical Modelling}

The machine under analysis is the fractional-slot PMSM assigned in Q2
(15~slots, 14~poles). All given parameters are listed in
Table~\ref{tab:q3_given}, and the intermediate geometric quantities
derived from them are listed in Table~\ref{tab:q3_calc}.

\begin{table}[h!]
\centering
\caption{Given design parameters.}
\label{tab:q3_given}
\begin{tabular}{llr}
\toprule
\textbf{Parameter} & \textbf{Symbol} & \textbf{Value} \\
\midrule
Number of phases              & $m$          & $3$ \\
Number of stator slots        & $Q$          & $15$ \\
Number of poles               & $N_m$        & $14$ \\
Slots per pole per phase      & $N_{spp}$    & $0.357$ \\
Stator outer radius           & $R_{so}$     & \SI{50}{mm} \\
Rotor/stator radius ratio     & $R_{ro}/R_{so}$ & $0.62$ \\
Motor axial length            & $L_s$        & \SI{100}{mm} \\
Air-gap length                & $g$          & \SI{1}{mm} \\
Radial magnet length          & $l_m$        & \SI{4}{mm} \\
Magnet relative permeability  & $\mu_r$      & $1.05$ \\
Magnet remanence              & $B_r$        & \SI{1.3}{T} \\
Magnet pole-arc ratio         & $\alpha_p$   & $0.889$ \\
Slot fill factor              & $k_{wc}$     & $0.60$ \\
Conductor current density     & $J$          & \SI{5}{A/mm^2} \\
Winding factor (fundamental)  & $k_{w1}$     & $0.957$ \\
Motor speed                   & $n$          & \SI{1500}{rpm} \\
DC bus voltage                & $V_{dc}$     & \SI{48}{V} \\
\bottomrule
\end{tabular}
\end{table}

Stator outer radius is given as $R_{so} = \SI{50}{mm}$, and the rotor-to-stator radius ratio is $0.62$, which allows us to calculate the rotor outer radius:
\[R_m = 0.62 \times R_{so} = 0.62 \times 50 = \mathbf{\SI{31.0}{mm}}\]
The stator inner radius is then:
\[R_{si} = R_m + g = 31.0 + 1.0 = \mathbf{\SI{32.0}{mm}}\]
The air-gap mean radius is the average of the rotor and stator radii:
\[R_{ag} = \frac{R_m + R_{si}}{2} = \frac{31.0 + 32.0}{2} = \mathbf{\SI{31.5}{mm}}\]

The teeth/slot width ratio is selected as $d = 0.62$, which allows to find outer radius of the slot, the height of the slot and the stator back-core thickness:
\[R_{slot,o} = \frac{R_{si}}{d} = \frac{32.0}{0.62} = \mathbf{\SI{42.67}{mm}}\]
\[h_s = R_{slot,o} - R_{si} = 42.67 - 32.0 = \mathbf{\SI{10.67}{mm}}\]
\[h_{core} = R_{so} - R_{slot,o} = 50.0 - 42.67 = \mathbf{\SI{7.33}{mm}}\]

Furthermore, the stator slot pitch is calculated as:
\[\tau_s = \frac{2\pi R_{ag}}{Q} = \frac{2\pi \times 31.5}{15} = \mathbf{\SI{13.40}{mm}}\]
If the teeth/slot width ratio is selected as $d = 0.62$, then the tooth body width and slot width are:
\[w_{th} = d \cdot \tau_s = 0.62 \times 13.40 = \mathbf{\SI{8.31}{mm}}\]
\[w_s = \tau_s - w_{th} = 13.40 - 8.31 = \mathbf{\SI{5.09}{mm}}\]
The slot area is:
\[A_{slot} = w_s \cdot h_s = 5.09 \times 10.67 = \mathbf{\SI{54.33}{mm^2}}\]
The copper area per slot is:
\[A_{cu} = k_{wc} \cdot A_{slot} = 0.60 \times 54.33 = \mathbf{\SI{32.60}{mm^2}}\]
The ampere-turns per slot is:
\[NI_{slot} = J \cdot A_{cu} = 5 \times 32.60 = \mathbf{\SI{163.0}{A}}\]
Finally, the electrical frequency at the rated speed is:
\[f_e = \frac{N_m \cdot n}{120} = \frac{14 \times 1500}{120} = \mathbf{\SI{175}{Hz}}\]

These calculated geometric parameters are summarized in Table~\ref{tab:q3_calc} below and used in the subsequent analytical calculations for the machine's performance characteristics.

\begin{table}[h!]
\centering
\caption{Calculated geometric parameters.}
\label{tab:q3_calc}
\begin{tabular}{llr}
\toprule
\textbf{Parameter} & \textbf{Symbol} & \textbf{Value} \\
\midrule
Rotor outer radius            & $R_m$          & \SI{31.0}{mm} \\
Stator inner radius           & $R_{si}$       & \SI{32.0}{mm} \\
Air-gap mean radius           & $R_{ag}$       & \SI{31.5}{mm} \\
Stator slot pitch             & $\tau_s$       & \SI{13.40}{mm} \\
Teeth/slot width ratio        & $d$            & $0.62$ \\
Tooth body width              & $w_{tb}$       & \SI{8.31}{mm} \\
Slot width                    & $w_s$          & \SI{5.09}{mm} \\
Slot outer radius             & $R_{slot,o}$   & \SI{42.67}{mm} \\
Slot height                   & $h_s$          & \SI{10.67}{mm} \\
Stator back-core thickness    & $w_{sy}$       & \SI{7.33}{mm} \\
Slot area                     & $A_{slot}$     & \SI{54.33}{mm^2} \\
Copper area per slot          & $A_{cu}$       & \SI{32.60}{mm^2} \\
Ampere-turns per slot         & $NI_{slot}$    & \SI{163.0}{A} \\
Stator slot opening           & $b_0$          & \SI{5.0}{mm} \\
Electrical frequency          & $f_e$          & \SI{175}{Hz} \\
\bottomrule
\end{tabular}
\end{table}


\subsection{Electrical Loading}

The specific electric loading $\bar{A}$ is the total RMS ampere-turns
per unit length of the air-gap circumference:

\begin{equation}
    \bar{A} = \frac{NI_{slot} \cdot Q}{\pi D_i}
            = \frac{163.0 \times 15}{\pi \times 0.063}
            = \mathbf{\SI{12{,}4}{kA/m}}
\end{equation}

\subsection{Carter's Coefficient and Effective Air-Gap}

Carter's coefficient corrects for the increase in effective reluctance
caused by stator slot openings. With selected slot opening $b_0 = \SI{5.0}{mm}$
and $g = \SI{1}{mm}$:

\begin{equation}
    \gamma = \frac{b_0/g}{5 + b_0/g}
           = \frac{5/1}{5 + 5}
           = 0.5
\end{equation}

\begin{equation}
    k_c = \frac{\tau_s}{\tau_s - \gamma\, b_0}
        = \frac{13.40}{13.40 - 0.5 \times 5.0}
        = \mathbf{1.23}
\end{equation}

Since the rotor has surface-mount magnets with no slots, the rotor
Carter's coefficient is unity ($k_{c2} = 1$). The effective mechanical
air-gap is then the physical air-gap multiplied by the stator Carter's coefficient:

\begin{equation}
    g' = k_c \cdot g = 1.23 \times 1 = \SI{1.23}{mm}
\end{equation}

\subsection{Effective Axial Length}

Fringing flux at the two axial ends extends the effective active length
by approximately one air-gap length per side:

\begin{equation}
    L_{eff} = L_s + 2g = 100 + 2\times 1 = \mathbf{\SI{102}{mm}}
\end{equation}

\subsection{Air-Gap Peak Flux Density}

Applying Ampere's law around a flux path that traverses the magnet and
the effective air-gap (cylindrical approximation, iron reluctance
neglected):

\begin{equation}
    B_{g,peak}
        = \frac{B_r \cdot l_m/\mu_r}{g' + l_m/\mu_r}
        = \frac{1.3 \times 4/1.05}{1.23 + 4/1.05}
        = \frac{4.952}{5.04}
        = \mathbf{\SI{0.983}{T}}
\end{equation}

\subsection{Magnetic Loading}

The specific magnetic loading is the average air-gap flux density over a
full pole pitch. Because the magnet covers only the fraction $\alpha_p$
of each pole, the remaining arc contributes zero flux:

\begin{equation}
    \bar{B} = \alpha_p \cdot B_{g,peak}
            = 0.889 \times 0.983
            = \mathbf{\SI{0.874}{T}}
\end{equation}

\subsection{Specific Machine Constant and Tangential Shear Stress}

The local tangential stress at the rotor surface is the product of the
local electric and magnetic loading. For unity power factor operation:

\begin{equation}
    \sigma_{tan} = \frac{\bar{A}\cdot\hat{B}_{g}}{\sqrt{2}}
                 = \frac{12{,}4 \times 0.983}{\sqrt{2}}
                 = \mathbf{\SI{8{,}59}{kPa}}
\end{equation}

This value is consistent with the expected range for a small
high-performance PMSM (\SIrange{5}{15}{kPa}). The specific machine
constant, combining winding design, electric loading, and magnetic
loading, is:

\begin{equation}
    C = \frac{\pi^2}{\sqrt{2}}\,k_{w1}\,\bar{A}\,\hat{B}_{g1}
      = \frac{\pi^2}{\sqrt{2}} \times 0.957 \times 12{,}4 \times 1.228
      = \mathbf{\SI{57{,}3}{kW\cdot s/m^3}}
\end{equation}

\subsection{Expected Torque}

The torque is related to the shear stress and the rotor volume:

\begin{equation}
    V_r = \pi R_m^2 L_{eff}
        = \pi \times (0.031)^2 \times 0.102
        = \SI{3.079e-4}{m^3}
\end{equation}

\begin{equation}
    T = 2\,\sigma_{tan}\,V_r
      = 2 \times 8{,}59 \times 3.079\times10^{-4}
      = \mathbf{\SI{5.29}{Nm}}
\end{equation}

\subsection{Flux Per Pole and Number of Turns}

The pole pitch and the flux per pole are:

\begin{equation}
    \tau_p = \frac{\pi D_i}{N_m}
           = \frac{\pi \times 0.063}{14}
           = \SI{14.1}{mm}
\end{equation}

\begin{equation}
    \Phi_p = B_{g,peak}\cdot\alpha_p\cdot\tau_p\cdot L_{eff}
           = 0.983 \times 0.889 \times 14.1\times10^{-3} \times 102\times10^{-3}
           = \mathbf{\SI{1.26}{mWb}}
\end{equation}

To select the number of turns, the target phase back-EMF is derived from
the \SI{48}{V} DC bus. Assuming maximum modulation index for the
three-phase inverter:

\begin{equation}
    E_{phase} = \frac{V_{dc}}{\sqrt{3}}
              \approx \SI{27.7}{V_{rms}}
\end{equation}


\begin{equation}
    N_s = \frac{E_{phase}}{\sqrt{2}\,\pi\,f_e\,k_{w1}\,\Phi_p}
        = \frac{27.7}{\sqrt{2}\times\pi\times175\times0.957\times1.260\times10^{-3}}
        \approx \mathbf{30\ \text{turns/phase}}
\end{equation}

\subsection{Tooth Flux Density Verification}

All air-gap flux entering one slot pitch must pass through the
corresponding tooth cross-section. The average flux density in the tooth body is:

\begin{equation}
    B_{tooth} = \frac{\bar{B} \cdot \tau_s}{w_{th}}
              = \frac{0.874 \times 13.40}{8.31}
              = \mathbf{\SI{1.409}{T}} \approx \SI{1.41}{T} \quad \checkmark
\end{equation}

\subsection{Results Summary}

\begin{table}[h!]
\centering
\caption{Summary of $Q3$ analytical modelling results.}
\label{tab:q3_summary}
\begin{tabular}{llr}
\toprule
\textbf{Quantity} & \textbf{Symbol} & \textbf{Value} \\
\midrule
Electrical loading          & $\bar{A}$        & \SI{12{,}4}{kA/m}          \\
Tooth flux density          & $B_{tooth}$      & \SI{1.41}{T}               \\
Carter's coefficient        & $k_c$            & $1.23$                     \\
Effective air-gap           & $g'$             & \SI{1.23}{mm}              \\
Total electromagnetic gap   & $g_{eff}$        & \SI{5.04}{mm}              \\
Effective axial length      & $L_{eff}$        & \SI{102}{mm}               \\
Peak air-gap flux density   & $B_{g,peak}$     & \SI{0.983}{T}              \\
Magnetic loading            & $\bar{B}$        & \SI{0.874}{T}              \\
Specific machine constant   & $C$              & \SI{57.3}{kW\cdot s/m^3}   \\
Tangential shear stress     & $\sigma_{tan}$   & \SI{8.59}{kPa}             \\
Expected torque             & $T$              & \SI{5.29}{Nm}              \\
Flux per pole               & $\Phi_p$         & \SI{1.26}{mWb}             \\
Number of turns per phase   & $N_s$            & $\approx 30$               \\
\bottomrule
\end{tabular}
\end{table}

All the calculations that are done in this question is implemented in an Excel spreadsheet named "analytical\_design.xlsx". The Excel file is provided in the project repository.



\section{Q4: 2D FEA Modeling}

To verify the analytical machine layout, a 2D FEA model was constructed using the ANSYS Maxwell software. The design flow was initiated in RMxprt as a \textit{Brushless Permanent Magnet DC Motor} template to properly and easily capture the surface-mounted permanent magnet synchronous motor (PMSM) topology. As the steel, M270-35A non-oriented thin electrical steel lamination is assigned for the stator and rotor cores. The magnets are defined as radially magnetized with a remanence of $B_r = 1.3\text{ T}$ and a relative permeability of $\mu_r = 1.05$. The magnet pole-arc ratio is set to $\alpha_p = 0.889$ to match the analytical design, and the slot fill factor is specified as $k_{wc} = 0.60$ for the winding regions.

The complete RMXprt analytical solution space was successfully validated, and an automated model export directly compiled the geometry, boundary planes, and mesh definitions into a 2D Transient Magnetic time-domain solver. In addition to the automated model, manual mesh refinement was performed in critical regions to ensure accurate field representation like air gap regions.

Additionally, to isolate the idealized performance characteristics of the PMSM under pure sinusoidal current excitation, several intermediate parameters were calculated and applied. The analytical slot loading was specified as $NI_{\text{slot}} = 163.0\text{ A}$ across 12 conductors per slot ($2 \times 6\text{ turns/coil}$). The rated rms phase current ($I_{\text{rms}}$) is derived as:
\begin{equation}
    I_{\text{rms}} = \frac{NI_{\text{slot}}}{\text{Conductors per Slot}} = \frac{163.0\text{ A}}{12} \approx 13.583\text{ A}_{\text{rms}}
\end{equation}
Converting this baseline into a peak current amplitude ($I_{\text{peak}}$) for a pure sinusoidal time function yields:
\begin{equation}
    I_{\text{peak}} = I_{\text{rms}} \cdot \sqrt{2} = 13.583 \times \sqrt{2} \approx 19.21\text{ A}
\end{equation}
The machine operates at a rated mechanical speed of $n = 1500\text{ rpm}$ with a 14-pole rotor, resulting in an electrical frequency of $f_e = 175\text{ Hz}$. A balanced three-phase stranded current excitation system was assigned directly to the phase windings via the following time-dependent formulas:
\begin{align}
    i_A(t) &= 19.21 \cdot \sin(2\pi \cdot 175 \cdot t) \\
    i_B(t) &= 19.21 \cdot \sin\left(2\pi \cdot 175 \cdot t - \frac{2\pi}{3}\right) \\
    i_C(t) &= 19.21 \cdot \sin\left(2\pi \cdot 175 \cdot t - \frac{4\pi}{3}\right)
\end{align}

Finalized 2D FEA model is shown in Figure~\ref{fig:fea_model}. The model was solved for a full mechanical cycle and results demanded in the project is presented in upcoming sections.

\begin{figure}[H]
    \centering
    \includegraphics[width=1\textwidth]{figures/sim_motor.png}
    \caption{Finalized 2D FEA model of the fractional-slot PMSM in ANSYS Maxwell.}
    \label{fig:fea_model}
\end{figure}


\subsection{Flux Density Distribution}

Figure~\ref{fig:flux_density_map} shows the flux density colour map of the full
machine cross-section at the aligned rotor position.

\begin{figure}[H]
    \centering
    \includegraphics[width=0.85\textwidth]{figures/flux_density_map.png}
    \caption{Flux density distribution in the machine cross-section.}
    \label{fig:flux_density_map}
\end{figure}

As can be seen from the figure, the flux density around the teeth for the aligned position is around \SI{1.3}{T}, which is smaller than the expected analytical value of \SI{1.41}{T}. This discrepancy can be attributed to the dicrepancies between the 2D FEA model and the analytical assumptions, such as the neglect of fringing effects and the idealized flux path in the analytical calculations. The FEA model captures more realistic flux distribution, including leakage and fringing, which can lead to a lower peak flux density in the teeth compared to the idealized analytical result. For the maximum flux density points, it can be said that the sharp edges of the teeth and the back of the magnets are the regions where the flux density is the highest, which is consistent with the expected behavior of a PMSM since these regions are where the flux lines are most concentrated.

\subsection{Air-Gap Flux Density Distribution}

The radial component of the flux density is plotted along a circular arc at the
middle of the air-gap ($r = R_{si} - g/2 = 32.5\;\text{mm}$) over one full
mechanical revolution, as shown in Figure~\ref{fig:airgap_flux}.

It can be seen that the air-gap flux density distribution is slightly non-sinusoidal, with sharp peaks corresponding to the magnet poles and deep valleys in between. The peak flux density in the air-gap is around \SI{0.95}{T}, which is close to the analytical prediction of \SI{0.983}{T}. The non-sinusoidal shape is expected due to the finite pole-arc ratio and the presence of slot openings, which cause flux concentration and leakage effects that deviate from the ideal sinusoidal distribution assumed in the analytical model. Also, due to saliency of the rotor, the flux density distribution is not perfectly sinudoidal.

\begin{figure}[H]
    \centering
    \includegraphics[width=1\textwidth]{figures/Air-gap Flux Density Distribution.png}
    \caption{Radial air-gap flux density at the mid air-gap radius over
             $360^\circ$ mechanical.}
    \label{fig:airgap_flux}
\end{figure}


\subsection{Leakage Flux Paths}

Figure~\ref{fig:leakage_flux} illustrates the flux line plot highlighting the
principal leakage paths present in the machine.

\begin{figure}[H]
    \centering
    \includegraphics[width=1\textwidth]{figures/leakage_flux.png}
    \caption{Flux line plot showing leakage paths.}
    \label{fig:leakage_flux}
\end{figure}

Two distinct leakage mechanisms are identified:
\begin{enumerate}
    \item \textbf{Magnet-to-magnet leakage:} Flux lines that leave one magnet
          pole and return directly to the adjacent pole without crossing the
          air-gap. This path is encouraged by the small inter-pole gap resulting
          from the $160^\circ$ electrical pole arc.
    \item \textbf{Magnet--tooth--magnet leakage:} Flux that crosses the air-gap
          into a stator tooth tip and then returns to an adjacent magnet rather
          than completing the main flux path through the back-core. This path is preferred since the local reluctance across the narrow slot opening and adjacent tooth tip is significantly lower than the total magnetic reluctance of the main structural path traveling all the way up the tooth body and through the stator back-core. 
\end{enumerate}

\subsection{Flux per Pole}

Each phase winding consists of 5 coils with 6 conductors each, giving a total
of $N_{phase} = 30$ series turns per phase. The peak flux linkage read from the
transient simulation is $\hat{\lambda} = 0.0302\;\text{Wb}$. Using the
fundamental winding factor $k_{w,1} = 0.957$ derived in Q2, the flux per pole is:

\[
    \Phi_p = \frac{\hat{\lambda}}{N_{phase}\cdot k_{w,1}}
           = \frac{0.0302}{30 \times 0.957}
           = \frac{0.0302}{28.71}
           = 1.052\times10^{-3}\;\text{Wb}
\]

The analytical flux per pole calculated in Q3 was $\Phi_{p,analytical} = 1.26\;\text{mWb}$.
The relative error between the two approaches is:
\[
    \varepsilon = \frac{|\Phi_{p,FEA} - \Phi_{p,analytical}|}{\Phi_{p,analytical}} \times 100
              = 16.7\;\%
\]

There exits a significant discrepancy between the analytical and FEA results for the flux per pole. This can be attributed to several geometric and physical effects that are not captured in the simplified analytical model. For example, the analytical model assumes an idealized flux path and neglects leakage flux, fringing effects, and the non-uniformity of the air-gap flux density. The FEA model captures these complex phenomena. Also, FEA model is tried to be arranged to be as close as possible to the analytical model, but there can still be some differences in the exact geometry, material properties, and boundary conditions that can lead to discrepancies in the results.

\subsection{Cogging Torque}

A transient simulation with zero winding current is run to find the cogging torque which is presented in Figure~\ref{fig:cogging_torque}. The cogging torque is the torque that arises due to the interaction between the permanent magnets and the stator teeth, even when no current is flowing in the windings.

\begin{figure}[H]
    \centering
    \includegraphics[width=1\textwidth]{figures/Cogging Torque.png}
    \caption{Cogging torque as a function of rotor angle.}
    \label{fig:cogging_torque}
\end{figure}



\section{Conclusion}

This project covered the complete winding design and analytical modelling workflow for two distinct electric machine configurations, verified through 2D finite element analysis. In Q1, the integral-slot winding analysis demonstrated that the 11/12 short-pitched configuration achieves a meaningful reduction in the 5th space harmonic winding factor at the cost of less than 1\% reduction in the fundamental winding factor. This confirms that short-pitching is an efficient strategy for harmonic suppression with negligible impact on torque and back-EMF capability. 

In Q2, the fractional-slot analysis of the 15-slot/14-pole machine yielded a fundamental winding factor of $k_{w,1} = 0.957$, which is high for a fractional-slot machine. The three-phase balance was verified analytically, confirming equal voltage magnitudes and exact $120^\circ$ phase displacement between phases. The analytical modelling in Q3 established the key electromagnetic loading parameters of the machine, providing reference values for comparison with FEA. 

The 2D FEA results in Q4 validated the analytical model. The flux density distribution confirmed that the teeth and back-core operate close to the design target of 1.4\,T. The air-gap flux density waveform and the flux linkage results were consistent with the analytical predictions with some deviations resulting from the simplified assumptions in the analytical model. The flux per pole extracted from the peak flux linkage of Phase A ($\hat{\lambda} = 0.0302$\,Wb, $N_{phase} = 30$ turns) yielded $\Phi_p = 1.052 \times 10^{-3}$\,Wb, in good agreement with the analytical calculations.


%%─── Appendix ────────────────────────────────────────────────────────────────
\appendix



\section{Project design parameters}
 Assigned design for this project is presented in Figure \ref{fig:design}.

\begin{figure}[H]
    \centering
    \includegraphics[width=0.7\textwidth]{figures/Design4_Ns15_Nm14.png}
    \caption{Assigned design visualization for the project.}
    \label{fig:design}
\end{figure}

\end{document}